\chapterhead{6}{What a Shadow Keeps}{Saying what the model keeps and drops}\label{ch6}

A shadow is a simple stand-in for a complicated thing. The shadow of a tree on the ground keeps the tree's outline and drops its colour, its texture and its leaves one by one. If you only had the shadow, you could still say a lot about the tree's height and the sun's position, but you could not say what colour the leaves were.

That is what a \textbf{model} does. A model keeps some things and drops others, on purpose. Every drawing in this book was a model of light.

\begin{center}
\sffamily\small\renewcommand{\arraystretch}{1.5}
\begin{tabular}{@{}p{3.6cm}p{3.7cm}p{3.7cm}@{}}
\toprule
\textbf{Model} & \textbf{What it keeps} & \textbf{What it drops}\\
\midrule
A straight line for light & the direction light goes & the colour, the brightness, the wave\\
A shadow & the object's outline and where the light is & colour, texture, inside\\
A dot for the sun & its direction and height & its size and brightness\\
A shadow clock & the hour from the tip's place & minutes, the season\\
A mirror rule & angle in equals angle out & bumpy surfaces\\
\bottomrule
\end{tabular}
\end{center}

\begin{worked}{Old measurements, reported}
Over two thousand years ago a Greek scholar named Eratosthenes is reported to have compared shadows at two places in Egypt on the same day to estimate how big the Earth is. The account comes from a later writer, Cleomedes, so historians treat the details with care. The idea is that a shadow's angle is a measurement that holds information about where on the globe you are standing.
\end{worked}

\begin{multix}[hmm]\multixbot[hmm]\tcblower
My shadow clock is a perfect model of the sun. It keeps everything about the sun. Is that right?
\end{multix}

\subhead{Five lines for a shadow report}
The report from the earlier books fits shadows, too.

\begin{tcolorbox}[enhanced,colback=white,colframe=ink!60,boxrule=1pt,arc=2mm,left=8pt,right=8pt,top=5pt,bottom=5pt,before skip=7pt,after skip=7pt]
\sffamily\small
\begin{enumerate}[leftmargin=1.6em,itemsep=0.2em]
\item \textbf{What I looked at.} The stick, the light, the place and the day.
\item \textbf{What I changed, measured and kept the same.}
\item \textbf{My readings.} Times and lengths or angles, all of them.
\item \textbf{My answer.} The pattern, with averages or errors where they apply.
\item \textbf{What I kept and what it hides.}
\end{enumerate}
\end{tcolorbox}

\clearpage

\begin{worked}{A sample report (practice numbers)}
\textbf{1. Looked at.} A 30~cm stick standing in clay on a taped sheet of paper, on a sunny day in the northern half of the world.\\
\textbf{2. Changed, measured, kept.} The time changed. I measured the shadow's length and direction. The stick, paper and place stayed the same.\\
\textbf{3. Readings.} Lengths at 8, 9, 10, 11, 12, 13, 14, 15 and 16:00: 62, 45, 34, 28, 26, 28, 35, 46 and 64~cm.\\
\textbf{4. Answer.} The shadow was shortest at 12:00, about 26~cm, and longest at the ends of the day. The shadow pointed west in the morning, north at midday and east in the afternoon.\\
\textbf{5. Kept and hidden.} I kept the shadow's length and direction. I set aside the sun's brightness, the day's weather and the season. A different day of the year may give different lengths.
\end{worked}

\begin{puzzle}{Safe sentences}
Mark each sentence S (supported by the shadow data), G (a guess about why) or B (too big for the data).
\begin{description}[leftmargin=1.6em,itemsep=0.3em,style=sameline]
\item[(a)] ``The shadow was 26~cm at 12:00.'' \wline[1cm]
\item[(b)] ``Shadows are shortest at noon everywhere on Earth.'' \wline[1cm]
\item[(c)] ``The shadow is short because the sun is high.'' \wline[1cm]
\end{description}
\end{puzzle}

\begin{puzzle}{Model or real thing?}
Write M if the sentence describes a model and R if it describes the real thing.
\begin{description}[leftmargin=1.6em,itemsep=0.3em,style=sameline]
\item[(a)] A drawing of a beam as a straight line. \wline[1cm]
\item[(b)] The sun in the sky. \wline[1cm]
\item[(c)] A paper clock with hour dots. \wline[1cm]
\end{description}
\end{puzzle}

\begin{puzzle}{What does it keep?}
A shadow clock keeps the direction of the sun's shadow. Write one thing it drops.\\
\rule{\linewidth}{0.5pt}
\end{puzzle}

\clearpage
\thispagestyle{fancy}\label{slop}
\subhead{Slop check}
Slop is hasty, gaudy work that forgot to take care. It can be made by a person or a machine, quickly or slowly. What matters is whether the work shows care for three things:
\begin{itemize}[leftmargin=1.6em]
\item \textbf{Salient:} Is the main point easy to see?
\item \textbf{Simple:} Is there anything that does nothing but get in the way?
\item \textbf{Relevant:} Does everything on the page help answer the question?
\end{itemize}

Here is a hasty poster about the shadow clock. Study it, then decide what is wrong.

\begin{center}
\begin{tikzpicture}[scale=0.97]
  \fill[yellow!40,draw=purple!80!black,line width=3pt] (0,0) rectangle (11.4,7.6);
  \node[star,star points=12,fill=red!60,draw=blue,line width=2pt,minimum size=2.2cm,font=\bfseries\scriptsize,rotate=10,text=yellow] at (1.4,6.2) {WOW!!!};
  \node[font=\itshape\bfseries\Large,text=teal,rotate=-4] at (6.4,6.8) {ShAdOwS aRe MaGiC};
  \node[font=\ttfamily\tiny,text=orange] at (6.4,6.15) {(sun maybe? who knows)};
  \foreach \x/\y/\c in {2.2/1.8/red,3.8/3.6/blue,5.4/2.4/green,7.0/4.2/orange,8.6/1.6/violet}{
    \fill[\c,draw=black] (\x,\y) circle (0.5);
  }
  \draw[thick,dashed,magenta] (2.2,1.8) -- (3.8,3.6) -- (5.4,2.4) -- (7.0,4.2) -- (8.6,1.6);
  \node[font=\tiny,text=black,rotate=60] at (9.8,3) {number number number};
  \foreach \x/\y in {1.2/1.5,1.6/4.2,10.2/5.6,0.8/3}{\node[star,star points=5,fill=yellow,draw=red,minimum size=0.6cm] at (\x,\y) {};}
  \node[font=\tiny,text=lime!50!black] at (5.4,0.45) {look at it!!!  best clock ever!!!  wow  wow};
\end{tikzpicture}
\end{center}

\begin{slopbox}
\textbf{(a)} What is the poster's main point? Can you tell?\\[3pt]\rule{\linewidth}{0.5pt}\\[6pt]
\textbf{(b)} Name two things that are not simple, because they only distract.\\[3pt]\rule{\linewidth}{0.5pt}\\[6pt]
\textbf{(c)} Name two things that are missing, so the dots cannot be read.\\[3pt]\rule{\linewidth}{0.5pt}\\[6pt]
\textbf{(d)} Which part is not relevant to the shadow clock at all?\\[3pt]\rule{\linewidth}{0.5pt}
\end{slopbox}

\clearpage
\thispagestyle{fancy}
\subhead{Now you try}
On a separate sheet, redraw a clear picture of the practice day: shadow length against time, from 8:00 to 16:00, with the numbers 62, 45, 34, 28, 26, 28, 35, 46 and 64 cm. Give it a title, label the axes with units, use one or two calm colours, and delete everything that does not help. Check it against the three words: salient, simple, relevant.

\begin{center}
\begin{tikzpicture}
  \draw[ink!70,line width=1pt,rounded corners=3pt] (0,0) rectangle (11.4,1.0);
  \node[font=\sffamily\small,anchor=west] at (0.2,0.5) {Title:};
  \draw[line width=1.3pt] (1,1.6) -- (1,9.6);\draw[line width=1.3pt] (1,1.6) -- (11,1.6);
  \foreach \y in {2.6,3.6,...,8.6}{\draw[gray!40] (1,\y) -- (11,\y);}
  \node[font=\sffamily\scriptsize,rotate=90] at (0.3,5.6) {label and unit here};
  \node[font=\sffamily\scriptsize] at (6,1.3) {label what the axis stands for};
\end{tikzpicture}
\end{center}

\clearpage
\thispagestyle{fancy}
\subhead{What did we keep, across the whole book?}
For each chapter, say one thing that was simplified, what was kept, and what was set aside. A sample is done for you.

\renewcommand{\arraystretch}{2.3}
\begin{center}\sffamily\scriptsize
\begin{tabularx}{\textwidth}{|>{\bfseries}p{2.1cm}|X|X|X|}
\hline
\textbf{Chapter} & \textbf{What we simplified} & \textbf{What we kept} & \textbf{What we set aside} \\ \hline
1 Light goes straight & light drawn as a line & the direction light goes & colour and brightness \\ \hline
2 Shadow size & & & \\ \hline
3 A stick in the sun & & & \\ \hline
4 A shadow clock & & & \\ \hline
5 Mirrors & & & \\ \hline
6 Shadow keeps & & & \\ \hline
\end{tabularx}
\end{center}

\begin{multix}[wow]\multixbot[wow]\tcblower
Everything you traced in this book was a simplification of something. A model is allowed to drop things, as long as you say which. Next time you see a long evening shadow, you will know what it is telling you.
\end{multix}
\clearpage
\thispagestyle{fancy}
\subhead{Mixed review}
\begin{puzzle}{Which is which?}
Write L (longest shadow), S (shortest shadow) or M (in between) for each time in the practice day: 8:00 \wline[1cm] 10:00 \wline[1cm] 12:00 \wline[1cm] 14:00 \wline[1cm] 16:00 \wline[1cm]
\end{puzzle}

\begin{puzzle}{Test detective}
(a) ``I measured the shadow of a stick on a hill at 9:00 and the shadow of a stick on a flat road at 3:00 and found the second was longer.'' What changed besides the time? \wline[4cm]\\
(b) ``I moved the lamp closer and the puppet away at once.'' What is wrong? \wline[4cm]
\end{puzzle}

\begin{puzzle}{Putting it together}
A 30~cm stick has a 34~cm shadow at 10:00. A building has a 340~cm shadow at the same time. How tall is the building? \wline[1.6cm]~cm \quad Hint: use the ratio.
\end{puzzle}

\begin{puzzle}{Short answer}
Name one thing a shadow clock keeps and one it drops.\\
\rule{\linewidth}{0.5pt}
\end{puzzle}
\clearpage
\thispagestyle{fancy}
\subhead{Shadow projects to try}
\begin{enumerate}[leftmargin=1.6em,itemsep=0.25em]
\item Measure the shadow of a tree or lamp post at three times. Use the ratio to work out its height.
\item Make a shadow puppet theatre and measure how big the puppets look from different distances.
\item Build a pinhole box and compare the picture size with the object size.
\item Draw the shadow of your house at 8:00, 12:00 and 16:00 on a sketch map.
\item Make a periscope from two small mirrors and a box.
\end{enumerate}
\renewcommand{\arraystretch}{2.4}
\begin{center}\sffamily\scriptsize
\begin{tabularx}{\textwidth}{|p{1.2cm}|X|X|X|}
\hline
\textbf{Project} & \textbf{What I changed} & \textbf{What I measured} & \textbf{What I kept the same} \\ \hline
 & & & \\ \hline
 & & & \\ \hline
 & & & \\ \hline
\end{tabularx}
\end{center}
\clearpage
\thispagestyle{fancy}
\subhead{My shadow diary}
Write down shadows you notice. Always include the time, where the sun was and how long or short the shadow was.
\renewcommand{\arraystretch}{2.4}
\begin{center}\sffamily\scriptsize
\begin{tabularx}{\textwidth}{|p{1.3cm}|p{1.3cm}|X|X|X|}
\hline
\textbf{Day} & \textbf{Time} & \textbf{What made the shadow} & \textbf{Which way it pointed} & \textbf{Length or size} \\ \hline
 & & & & \\ \hline
 & & & & \\ \hline
 & & & & \\ \hline
 & & & & \\ \hline
 & & & & \\ \hline
 & & & & \\ \hline
 & & & & \\ \hline
 & & & & \\ \hline
\end{tabularx}
\end{center}
\clearpage
\thispagestyle{fancy}
\subhead{Questions I still have}
Write a question about light or shadows that this book made you wonder about. Then write how you might find out.
\par\vspace{1em}
\foreach \i in {1,...,12}{\noindent\rule{\linewidth}{0.4pt}\par\vspace{20pt}}
\clearpage
